\section{Model Quantization}

\subsection{Khái niệm Quantization}

Quantization ánh xạ giá trị số thực có độ chính xác cao sang biểu diễn có số bit
thấp hơn. Với lượng tử tuyến tính, một giá trị $x$ được ánh xạ thành số nguyên
$q$ thông qua scale $s$ và zero-point $z$:
\begin{equation}
  q=\operatorname{clip}\left(\operatorname{round}\left(\frac{x}{s}\right)+z\right).
\end{equation}
Việc giảm số bit làm giảm dung lượng và có thể tăng thông lượng, nhưng gây sai số
lượng tử cần được kiểm soát.

\subsection{FP32, FP16 và INT8}

FP32 là biểu diễn phổ biến khi huấn luyện và được dùng làm chuẩn chất lượng.
FP16 giảm một nửa dung lượng mỗi phần tử, phù hợp với Tensor Core và thường giữ
accuracy gần FP32. INT8 dùng số nguyên 8 bit, có tiềm năng giảm bộ nhớ và tăng
tốc mạnh hơn nhưng nhạy với dải giá trị activation và yêu cầu calibration hoặc
QAT.

\subsection{Post-Training Quantization}

Post-Training Quantization thực hiện lượng tử sau khi mô hình đã huấn luyện. PTQ
không yêu cầu huấn luyện lại hoặc chỉ cần calibration dataset để ước lượng dải
activation. Ưu điểm là quy trình nhanh; hạn chế là accuracy có thể giảm với tầng
nhạy cảm, dữ liệu lệch hoặc mô hình nhỏ.

\subsection{Quantization-Aware Training}

QAT mô phỏng hiệu ứng lượng tử trong quá trình huấn luyện để trọng số thích nghi
với sai số làm tròn và clipping \cite{jacob2018}. Forward pass sử dụng các nút
fake quantization, còn gradient được xấp xỉ để cập nhật trọng số thực. QAT tốn
thêm thời gian huấn luyện nhưng thường duy trì accuracy INT8 tốt hơn PTQ.

\subsection{Fake Quantization}

Fake quantization lượng tử và giải lượng tử tensor trong forward pass nhưng vẫn
lưu trọng số ở dạng số thực để tối ưu. Các observer thu thập thống kê min--max
hoặc histogram nhằm xác định scale. Khi export, các thông tin này được chuyển
thành toán tử lượng tử hoặc metadata phục vụ TensorRT.

\subsection{So sánh PTQ và QAT}

\begin{table}[H]
\centering
\caption{So sánh PTQ và QAT}
\label{tab:ptq-qat}
\begin{tabularx}{\textwidth}{L{3.2cm}YY}
\toprule
\textbf{Tiêu chí} & \textbf{PTQ} & \textbf{QAT}\\
\midrule
Dữ liệu cần thiết & Tập calibration đại diện & Dữ liệu huấn luyện/validation\\
Chi phí thực hiện & Thấp, không huấn luyện lại & Cao hơn do fine-tuning có fake quant\\
Khả năng giữ accuracy & Phụ thuộc độ nhạy mô hình & Thường tốt hơn với INT8\\
Độ phức tạp pipeline & Đơn giản & Cần cấu hình và export tương thích\\
Vai trò trong đồ án & Baseline INT8 tham chiếu & Phương pháp tối ưu chính\\
\bottomrule
\end{tabularx}
\end{table}